\section{Background}\label{sec:background}

\method\ uses graph theory to represent both the street network and the sewage collector network.
The urban mesh is modeled as a graph $G=(V,E)$, where $V$ is the set of relevant points in urban space and $E$ is the set of connections between them~\cite{Boeing2017}.
Vertices correspond to intersections, dead-ends, curves, intermediate points, collection points, or topographic points of interest.
Edges correspond to street segments or collector segments, carrying attributes such as length, geometry, road type, and elevation.

This representation enables classical graph properties to be used in interpreting the urban mesh~\cite{Turan2019}.
The degree of a vertex indicates how many connections reach a point, helping to identify crossings and bifurcations.
Paths represent possible sequences for flow conduction.
Connectivity indicates whether all regions of the mesh remain reachable from a discharge or collection point, which is a fundamental requirement for ensuring full network coverage, avoiding isolated segments, and identifying points where direction changes or maintenance requirements imply physical structures such as manholes.

The input street network is initially treated as an undirected graph, since streets provide geometric connections between adjacent points without defining a direction of sewage flow.
After topographic analysis, the resulting collector network is represented as a directed graph, where each arc indicates a preferred flow direction~\cite{Haydar2026}.
Gravity is a fundamental physical constraint in sanitary sewer systems: flow must be conducted from higher elevations to lower elevations without pumping wherever possible~\cite{Turan2019}.
Topography therefore directly influences path selection, edge orientation, and the interpretation of high points, low points, and grade breaks.

Since automatically extracted street data from sources may contain redundant intermediate nodes, the model also requires graph simplification~\cite{Boeing2017}.
The challenge is to reduce topological redundancy without discarding technically relevant elements, such as corners, sharp curves, collection points, and nodes constrained by maximum spacing requirements between manholes.

